\section{RによるWithinデザインの分散分析}
\label{b21_WithinDesign_onR}

\subsection{授業内容}
\subsubsection{概要}
前回学んだ群内計画(Within-subjects Design)の理論的理解をもとに，本コマではRを用いた群内計画の実践的な分析方法について学ぶ。特に，anovakunという関数群を活用し，群内計画データの構造化，モデルの構築，分散分析の実施，結果の解釈までの一連のプロセスを習得する。群内計画特有の個人差の分離，効果の検出，交互作用の解釈についてRの出力結果から理解を深める。また，分析結果の適切な可視化と報告方法についても学び，心理統計学の実践的スキルを向上させる。

\subsubsection{コマ主題細目}
\begin{description}

	\item[anovakunによる一要因群内計画の分析] anovakunパッケージを用いた一要因群内計画の分析方法を習得する。ANOVA関数と群内計画を指定するための構文（参加者ID変数の指定方法など），モデル式の記述法を理解する。得られた分散分析表の読み方と解釈，特に個人差（参加者間変動）が明示的に分離されていることを確認する。効果量（$\eta^2$，$\omega^2$など）の算出方法と，分析結果の適切な報告形式についても習得する。
	      \begin{flushright}
		      $\to$ \textcite{kosugi2019}のP.152--157.
	      \end{flushright}

	\item[anovakunによる多要因群内計画の分析] 二要因以上の群内計画の分析方法について学ぶ。複数の実験要因がすべて群内要因である場合のモデル式の記述法，交互作用の指定方法，結果の解釈方法を理解する。特に，交互作用が有意であった場合の単純主効果の検定方法についてanovakunでの実施方法を習得する。また，多要因群内計画における効果の分解と誤差項の選択，適切なF値の算出方法についても学ぶ。実際の心理学実験データを用いた演習を通じて，多要因群内計画の分析スキルを向上させる。
	      \begin{flushright}
		      $\to$ \textcite{kosugi2018}のP.125--132.
	      \end{flushright}

	\item[混合計画の場合] 群内要因と群間要因が混在する混合計画の分析方法について学ぶ。anovakunでのモデル式の記述法，交互作用の指定方法，結果の解釈方法を理解する。特に，混合計画における個人差の分離と効果量の算出方法についても学ぶ。実際の心理学実験データを用いた演習を通じて，混合計画の分析スキルを向上させる。

\end{description}
\subsubsection{キーワード}
\begin{itemize}
	\item anovakunパッケージによる群内計画分析
	\item 一要因および多要因群内計画の分散分析
	\item 球状性の検定
	\item 混合計画の分析
	\item 効果量の算出と解釈
\end{itemize}
\subsection{授業情報}
\paragraph{コマの展開方法}
講義とR実習
\paragraph{標準シラバスにおける位置づけ}
科目番号4；心理学研究法；2.データを用いた実証的な思考方法；C データの統計的記述\\
科目番号5；心理学統計法；1.心理学で用いられる統計手法；G 推測統計:代表値と散布度をめぐって(1)
\subsubsection{予習・復習課題}
\paragraph{予習}
前回学習した群内計画(Within Design)の理論的内容を復習しておくこと。特に，群内計画における個人差の分離の意義，平方和の分解，適切なF値の算出方法について理解を確認しておくとよい。また，Rの基本操作，データフレームの扱い方，基本的な関数の使用方法についても再確認しておくこと。

\paragraph{復習}
授業で学んだ内容を定着させるために，以下の点について理解を深めておくこと：

\begin{enumerate}
\item 群内計画データの適切な構造化方法を理解し，wide形式とlong形式の相互変換ができるようにする
\item anovakunパッケージを用いた一要因群内計画の分析を独力で実施し，結果を正しく解釈できるようにする
\item 多要因群内計画の分析方法を理解し，交互作用の解釈と単純主効果の検定ができるようにする
\item 群内計画の分析結果を適切に可視化し，APA形式に従った報告ができるようにする
\item 実際の研究データや公開データを用いて，群内計画の分析を練習し，スキルを向上させる
\end{enumerate}

授業で扱った分析スクリプトを，異なるデータセットに適用してみることで理解を深めるとよい。公開されている心理実験データや教科書に掲載されているデータを用いて実習してみるとよい。また，各自の研究計画において群内計画が適用可能かどうかを検討し，必要に応じて分析計画を立ててみるとよい。なお，anovakunの開発者のWebサイトやGitHubページで提供されている追加の解説やチュートリアルも参考になる。